

```latex
\documentclass{article}
\usepackage{pathfinder-note}

\title{Adversarial Review for Market Convergence: A Cross-Paper Analysis}

\begin{document}

\maketitle

\section*{The Pair}

Q examines LLM agents in double-auction markets, finding slower or absent convergence to equilibrium compared to human traders, with substantial heterogeneity across model families and market roles. CoT analysis reveals trading decisions correlate with a shift from strategic considerations toward urgency~\ledger{1, 3}.

P introduces Adversarial Review (AR), a three-agent protocol where a reviewer evaluates code and a critic audits through structured disagreement before the main agent edits. AR outperforms larger multi-agent baselines on code-review tasks, specifically addressing false-consensus failure modes~\ledger{1, 3}.

\section*{Connexion Pursued}

Initial hypothesis: P's structured disagreement protocol could improve Q's market convergence. The proposed mechanism was that AR-style reviewer-critic friction might interrupt premature trade commitments, forcing evidence-grounded justification before execution~\ledger{1, 3}.

\section*{Supported Findings}

The peers established:

\begin{itemize}
  \item Q's core pathology is \emph{heterogeneity} (agents failing to coordinate), not false-consensus. These are opposite failure modes~\ledger{2, 5}.
  \item Q's CoT finding (strategic\(\rightarrow\)urgency shift) operates at the \emph{individual} agent level, while heterogeneity is a \emph{market-level} phenomenon~\ledger{6}.
  \item AR is deliberative and collaborative; double-auction trading is competitive and time-sensitive. Adding disagreement overhead could worsen convergence by introducing latency~\ledger{4}.
  \item A narrow salvageable connection exists only if: (1) urgency-driven trades are insufficiently justified at the individual level, (2) self-critique reduces decision noise, and (3) this improves market-level coordination~\ledger{6}.
\end{itemize}

\section*{Objections}

\begin{enumerate}
  \item \textbf{Pathology mismatch}: Q shows heterogeneity (insufficient agreement); P addresses false-consensus (excessive agreement without evidence). These are opposite pathologies~\ledger{2, 5}.
  \item \textbf{Domain mismatch}: Code review tolerates latency; double auctions are time-sensitive. AR friction could harm rather than help~\ledger{4}.
  \item \textbf{Role mapping unclear}: P's reviewer/critic roles have no obvious analogue in trading. What counts as ``evidence'' for a trade remains undefined~\ledger{6}.
\end{enumerate}

\section*{Unresolved Issues}

Evidence is thin on:

\begin{itemize}
  \item Whether urgency-driven trades are the \emph{cause} of non-convergence or a symptom of deeper misunderstanding of market dynamics~\ledger{4}.
  \item Whether heterogeneity is structural (different models/roles inherently incompatible) or a result of correctable decision-making patterns~\ledger{5}.
  \item What a trading-adapted AR protocol would look like: timing, evidence standards, and whether self-critique (vs. multi-agent) suffices~\ledger{6}.
\end{itemize}

\section*{Smallest Next Step}

Examine Q's full CoT traces to determine whether individual agents exhibit premature commitment patterns that resemble false-consensus at the decision level (before trade execution), even while market-level outcomes show heterogeneity. This would establish whether AR-style interruption is mechanistically plausible.

\end{document}
```